\chapter{Conclusiones}
\label{chap:conclusiones}

A lo largo de este trabajo se recorrió el ciclo completo de desarrollo de filtros activos: el cálculo a partir de
una especificación, la normalización de los componentes, la simulación con elementos ideales y con el modelo del
operacional real, el diseño del circuito impreso y, finalmente, la medición en laboratorio. Ese recorrido se aplicó a
cuatro circuitos: el módulo de acondicionamiento de la señal de campo de $\SI{4}{\milli\ampere}$ a
$\SI{20}{\milli\ampere}$, implementado y caracterizado en forma estática y dinámica, los filtros pasa bajos y pasa
altos de cuarto orden, montados sobre una misma placa y relevados en frecuencia, y el filtro pasa banda de segundo
orden, que se diseñó y simuló y cuya implementación queda para el trabajo de filtros digitales. A la luz de los
resultados obtenidos, es posible evaluar el cumplimiento de los objetivos planteados en la
sección~\ref{sec:objetivos}.

En el módulo de señal de campo se cumplió el objetivo de comprender el funcionamiento del receptor diferencial y del
convertidor tensión-corriente: la caracterización estática mostró que el módulo reproduce a su salida la corriente de
entrada con transferencia unitaria y lineal en todo el rango de $\SI{4}{\milli\ampere}$ a $\SI{20}{\milli\ampere}$.
La caracterización dinámica verificó el objetivo de filtrado: la frecuencia de corte resultó la de diseño y la
componente de red de $\SI{50}{\hertz}$ se atenúa en la medida esperada.

En cuanto a los tres filtros asignados, las simulaciones con valores calculados y operacional ideal reprodujeron
exactamente las especificaciones de la tabla~\ref{tab:asignacion}, lo que confirma que las técnicas de cálculo vistas
en clase (aproximación de Butterworth, transformaciones de frecuencia y dimensionamiento de las celdas Sallen-Key por
proporcionalidad de componentes) se aplicaron correctamente. Con valores comerciales y el TL084, los filtros pasa
bajos y pasa altos siguieron cumpliendo la atenuación máxima en $f_p$. El pasa banda, en cambio, no cumple la
especificación con el operacional elegido: la atenuación dentro de la banda supera ampliamente el límite de
$A_{\max}$. La causa es que la única celda opera cerca de su ganancia de oscilación
($G / G_{osc} = 0{,}90$), lo que la hace muy sensible: el pequeño retardo de fase que introduce el operacional en la
frecuencia central basta para desplazar los polos y correr la banda de paso fuera de la especificada.

El filtro pasa bajos implementado cumplió todas sus especificaciones en el rango en que la medición fue posible: la
ganancia de $\SI{4}{\volt/\volt}$ en la banda de paso, la atenuación en $f_p$ menor que $A_{\max}$ y la frecuencia de
corte de diseño. Su respuesta medida coincide con la simulada, ya que su banda de paso termina muy por debajo del
ancho de banda a lazo cerrado de cada celda y el TL084 se comporta prácticamente como un operacional ideal.

El pasa altos, por el contrario, puso en evidencia las limitaciones del operacional real, tal como había anticipado
la simulación. La ganancia máxima medida coincide con la simulada con el TL084, pero es menor que los $\SI{12.25}{\volt/\volt}$ calculados y que los $\SI{12.5}{\volt/\volt}$ especificados. La diferencia con
el cálculo se debe al producto ganancia-ancho de banda del TL084: con una ganancia de $3{,}5$ por celda, el ancho de
banda a lazo cerrado es de apenas unas 27 veces la frecuencia de corte, y el desfasaje del operacional altera además
los polos dentro de la banda de paso. A esto se suma que en el diseño se adoptó $G = 3{,}5$ por celda, es decir
$\SI{12.25}{\volt/\volt}$ en lugar de $\SI{12.5}{\volt/\volt}$, por lo que la ganancia especificada no se habría
alcanzado aun con un operacional ideal. El objetivo de ganancia, por lo tanto, no se cumple. La forma de la respuesta,
en cambio, sí satisface la especificación: tomando como referencia la ganancia real de la banda de paso, la
atenuación en $f_p$ es menor que $A_{\max}$ y la frecuencia de corte coincide con la de diseño. Por encima de la banda
de paso apareció además el corte superior previsto por la simulación, con lo que el filtro real se comporta como un
pasa banda muy ancho.

La FFT de la salida durante el barrido de frecuencia mostró en ambos filtros el perfil característico buscado: en el
pasa bajos, un nivel plano en la banda de paso que cae por encima del corte, y en el pasa altos, la banda de rechazo
en baja frecuencia seguida de la subida hasta el nivel de la banda de paso. Esto confirma de forma cualitativa el tipo
de respuesta de cada placa. La lectura de la frecuencia de corte sobre esa FFT dio, en cambio, en los dos filtros un
valor menor que el de diseño, con un corrimiento más marcado en el pasa bajos que en el pasa altos. Este resultado no coincide con
el relevamiento punto a punto, que ubicó la caída de la respuesta en la frecuencia de corte de diseño, y se atribuye
al método más que a los filtros. El perfil que muestra la FFT durante el barrido no es solo el módulo de la
transferencia: incluye también la forma en que el barrido reparte su energía en frecuencia y el ensanchamiento propio
de superponer muchos espectros, por lo que su envolvente no es plana aun en la banda de paso. A esto se suma que la
lectura con cursores es gruesa, al punto que la separación entre los cursores de amplitud no quedó exactamente en el
valor buscado. El método resulta útil para obtener rápidamente una estimación de la frecuencia de corte, pero para
verificar la especificación se toma como válido el relevamiento punto a punto.

Las mediciones no permitieron verificar, en ninguno de los dos filtros, la atenuación en el extremo de la banda de
rechazo: en $10 f_p$ del pasa bajos y en $f_p/10$ del pasa altos se leyeron tensiones muy superiores a las
esperadas. Esas lecturas corresponden al piso de ruido e interferencia del conjunto punta y osciloscopio y no a la salida de los
filtros. La causa es que todos los puntos se midieron con una entrada de $\SI{1}{\volt_{\text{pp}}}$, sin aplicar el
aumento de la señal ni el promediado previstos en las condiciones de ensayo.

En conjunto, los objetivos generales del trabajo se cumplieron: se estudiaron y aplicaron las técnicas de análisis y
diseño de filtros activos y se implementaron tres circuitos cuyo comportamiento medido coincide con el calculado y el
simulado. La enseñanza principal es que el operacional real deja de ser transparente cuando la banda de interés se
acerca a su ancho de banda a lazo cerrado, como en el pasa altos, o cuando la celda trabaja con alta sensibilidad,
como en el pasa banda; en esos casos la simulación con el macromodelo resultó una herramienta confiable para
anticipar lo que luego se midió, mientras que el cálculo ideal no alcanza.

Quedan abiertas algunas líneas de trabajo que surgen de estas limitaciones. La banda de rechazo de los filtros pasa
bajos y pasa altos debería relevarse nuevamente con una entrada de $\SI{10}{\volt_{\text{pp}}}$ y con promediado o
FFT, y los puntos de alta frecuencia del pasa altos con $\SI{0.5}{\volt_{\text{pp}}}$, para descartar la influencia
del \emph{slew rate}. Para alcanzar la ganancia especificada del pasa altos sería
necesario recalcular la red $R_a$--$R_b$ para $\SI{12.5}{\volt/\volt}$ y emplear un operacional de mayor producto
ganancia-ancho de banda, o bien predistorsionar el diseño para compensar su efecto; la misma estrategia permitiría
que el pasa banda cumpla la especificación de $A_{\max}$. Finalmente, la implementación del pasa banda en el trabajo
práctico de filtros digitales ofrecerá la oportunidad de comparar una realización analógica y una digital de una misma
especificación, en particular frente a la sensibilidad que este trabajo puso de manifiesto.
